\section{模型假设}
为使模型既可计算又不过度偏离题意，本文作如下假设。

\begin{enumerate}[itemindent=2em]
  \item 板凳在运动中视为刚体，单节板凳的几何尺寸在整个过程中保持不变。
  \item 每节板凳的两个把手中心位置固定，且相邻节之间仅通过把手连接。
  \item 盘入和盘出阶段的路径均可视为理想平面曲线，忽略弹性、摩擦和摆动等细节。
  \item 除题目允许的相邻连接外，任意两节非相邻板凳一旦实体相交即视为发生碰撞。
  \item 数值计算中的极小残差仅来自离散化与求根误差，不改变几何判定的结论。
\end{enumerate}

上述假设中，第4条是碰撞判定的核心，第5条用于说明数值结果的适用边界。它们共同保证后续结论具有明确的几何意义。
